% TOP_PROBLEM: {"id": 4600022, "title": "Stable cellular-automaton limit sets", "category_id": 11, "category": {"id": 11, "name": "dynamical_systems", "display_name": "Dynamical Systems", "description": "Problems about long-term behavior of deterministic systems, Hamiltonian dynamics, and periodic orbits.", "slug": "dynamical-systems", "order_index": 11, "created_at": "2026-07-31T15:26:25.671Z"}, "status": "open", "problem_number": "AMR-045-0022", "set_id": 13}
% FIRST_POSTED: 2026-10-01
% ENTRY_KIND: statement_only
% UnsolvedMath ID 4600022; source catalogue code AMR-045-0022
% !TeX program = lualatex
% Definition and sources only; this file is not an LLM solution attempt.
\documentclass[11pt,a4paper]{article}
\usepackage[margin=25mm]{geometry}
\usepackage{fontspec}
\setmainfont{lmroman10-regular.otf}[BoldFont=lmroman10-bold.otf,
 ItalicFont=lmroman10-italic.otf,BoldItalicFont=lmroman10-bolditalic.otf]
\usepackage{amsmath,amssymb,mathtools}
\usepackage{unicode-math}
\setmathfont{latinmodern-math.otf}
\AtBeginDocument{\let\setminus\smallsetminus}
\usepackage{xurl,hyperref}
\hypersetup{colorlinks=true,urlcolor=blue}
\setcounter{secnumdepth}{0}
\setlength{\parindent}{0pt}
\setlength{\parskip}{5pt}
\setlength{\emergencystretch}{3em}
\begin{document}
\section*{Stable cellular-automaton limit sets}\label{4600022}
{\small UnsolvedMath ID 4600022 · AMR-045-0022 · Dynamical Systems · AMR Open Problem Lists}
\subsection{Definitions and mathematical statement}
Let $A$ be a finite nonempty alphabet. A one-dimensional cellular
automaton is a continuous map $F:A^{\mathbb Z}\to A^{\mathbb Z}$
commuting with the left shift. Equivalently, some radius $r\ge0$ and
local rule $\varphi:A^{2r+1}\to A$ satisfy
$(Fx)_j=\varphi(x_{j-r},\ldots,x_{j+r})$ at every site.
The decreasing sequence of compact images defines its limit set
\[
       \Omega(F)=\bigcap_{n\ge0}F^n(A^{\mathbb Z}).
\]
The limit set is stable if there is a finite $k\ge0$ with
$F^k(A^{\mathbb Z})=F^{k+1}(A^{\mathbb Z})$; then this common image
equals $\Omega(F)$.

Characterize the subshifts which arise as stable limit sets of such
automata. One formulation is to give necessary and sufficient conditions
on $Y\subset A^{\mathbb Z}$ for the existence of a finite-window rule
$F$ and a finite stabilization time $k$ with $F^k(A^{\mathbb Z})=Y$.
The corresponding formulation up to shift conjugacy allows a finite
recoding of the target alphabet.

Every finite-time image is a sofic shift, since it is a factor of a
full shift; thus the problem lies within the finite labelled-graph
class. The classification must identify which such spaces admit a
surjective dynamics induced by a local rule on the surrounding full
shift, with all configurations reaching the image in a uniform finite
time. Pointwise eventual approach alone is not stability.
\subsection{Short English statement}
Which symbolic spaces are reached, in a uniformly bounded number of steps, as the final image of a cellular automaton?
\subsection{Sources}
\begin{itemize}
\item[S1] ULAM.AI and the UnsolvedMath Contributors, supplied problems.json, record 4600022 (AMR-045-0022), AMR Open Problem Lists. Catalogue identity and source transcription; CC BY 4.0.
\url{https://huggingface.co/datasets/ulamai/UnsolvedMath}.
\item[S2] Boyle - Open Problems in Symbolic Dynamics (2008), Problem/Question/Conjecture 16.1. Original source indexed by AMR-045-0022.
\url{https://www.math.umd.edu/~mboyle/papers/openfinalsub3nov2007.pdf}.
\end{itemize}
\end{document}
